% !PRINTED 643
\weberpage{643}

%<<CONT
sind, gehen also, wenn für die $x$ die Substitution $x_h = \psi(\eta_h)$ gemacht wird, in Functionen der Variablen $a$ über, die auch nicht identisch verschwinden können. Nach dem Satze \S.\,143, 1. kann man also für die Grössen $a_0, a_1 \ldots a_{n-1}$ solche rationale Zahlen setzen, dass die Functionen $(\varepsilon, x)$ und die Differenzen $f_\alpha - f_\beta$, die dann also in $(\varepsilon, \xi)$ und in $k_\alpha - k_\beta$ übergehen, von Null verschieden werden, und dass auch die $\xi_h$ von einander verschieden bleiben, was zu beweisen war.

Nachdem also dies festgestellt ist, führen wir das System der Variablen $x_0, x_1 \ldots x_{n-1}$ und ein davon abhängiges System von Variablen $y$

\weq{5}{y_0 = \chi(x_0), \quad y_1 = \chi(x_1), \ldots y_{n-1} = \chi(x_{n-1})}

ein, worin $\chi$ die Function (3) bedeutet, und setzen

\weq{6}{\dfrac{(\varepsilon^{r_\nu}, y)}{(\varepsilon^{r_\nu}, x)} = \Theta_\nu.}

Wird auf die Indices der $x$ irgend eine Permutation angewandt, so erleiden die Indices der $y$ die gleiche Permutation. Wenn wir also auf (6) die cyklische Permutation $s$ anwenden, so ändert sich $\Theta_\nu$ nicht (nach \S.\,184, 4.).

Wenn wir also $\Theta_\nu$ in der Normalform des \S.\,183 darstellen, so sind seine Coëfficienten cyklische Functionen von $x_0, x_1 \ldots x_{n-1}$.

Wendet man auf $\Theta_\nu$ die Permutation $t^{-1} = (h, g^{-1}h)$ an, so geht $\Theta_\nu$ nach \S.\,184, 4. in $\Theta_{\nu+1}$ über, oder die $\Theta_\nu$ erleiden eine cyklische Permutation.

Denselben Erfolg hat aber auch die Substitution $\sigma = (\varepsilon, \varepsilon^g)$, und also hat das System der Functionen $\Theta_0, \Theta_1 \ldots \Theta_{n-2}$ die Eigenschaften $\alpha)\ \beta)\ \gamma)\ \delta)$ (\S.\,184).

Nach dem am Ende des \S.\,184 bewiesenen Satze geht also, wenn wir die $x$ durch die $\xi$ und folglich die $y$ durch die $\eta$ ersetzen, $\Theta_\nu$ in eine rationale Function von $k_\nu$

\weq{7}{Q_\nu = Q(k_\nu)}

über, und aus (6) ergiebt sich

\weq{8}{(\varepsilon^{r_\nu}, \eta) = Q_\nu (\varepsilon^{r_\nu}, \xi).}

Wenn man hierin für $(\varepsilon^{r_\nu}, \xi)$ den Ausdruck (5), \S.\,185 substituirt, so erhält man für $(\varepsilon^{r_\nu}, \eta)$ einen Ausdruck von ganz derselben Form, nur mit dem Unterschiede, dass an Stelle von $K_\nu$ getreten ist $Q_\nu K_\nu$. Die Functionen $Q_\nu K_\nu$ haben im Wesentlichen dieselben Eigenschaften, wie die Function $K_\nu$, nur dass sie auch
%CONT>>
